% Source-ready standalone technical section. Full proofs, no external inputs.
% 5 November 2026. Computational companions are in this directory.
\documentclass[11pt]{article}
\usepackage[T1]{fontenc}
\usepackage[margin=1in]{geometry}
\usepackage{amsmath,amssymb,amsthm}
\usepackage[hidelinks]{hyperref}
\newtheorem{theorem}{Theorem}[section]
\newtheorem{proposition}[theorem]{Proposition}
\newtheorem{lemma}[theorem]{Lemma}
\newtheorem{definition}[theorem]{Definition}
\theoremstyle{remark}\newtheorem{remark}[theorem]{Remark}
\newcommand{\Pow}{\mathcal P}
\newcommand{\Gen}{\operatorname{Gen}}
\newcommand{\Meet}{\operatorname{Meet}}
\newcommand{\cl}{\operatorname{cl}}
\newcommand{\im}{\operatorname{im}}
\title{Operationally generated incidence under changing experimental interfaces}
\author{Source-ready mathematical development}
\date{5 November 2026}
\begin{document}\maketitle

\section{Scope and independent source data}
The question is not merely how three invariants of one table differ.
An interaction can make a new experiment executable, and a new operation
can turn an old observation into a new experiment by precomposition.
We derive the resulting comparisons from these source operations. There
are two comparison directions and an essential distinction between
strict, lax, and partial transport.

Let $\mathcal C$ be a small category of admitted native programs and
$X:\mathcal C\to\mathsf{Set}$ their deterministic state semantics.
Objects are experimental interfaces. At interface $i$, supply a set
$O_i$ of labelled seed observations, with their actual Boolean response
columns $C_e\subseteq X_i$. A label includes a native observation
witness; a morphism of $\mathcal C$ includes its native program witness.
The functor $X$ need not be faithful: two different native programs may
have the same state transformation.

This is a relative source model, not a theorem that finitely many
observations determine all physical laws or all possible programs.
Actual residual stock, damaged specimens, and installed devices belong
in the source state when they affect subsequent execution. A partial
command can be represented on a \emph{certified domain interface} on
which it is total; this does not turn unsafe disabledness into an
observable outcome. A measurement may have a state-changing native
implementation even though its final Boolean response alone supplies
an incidence column. In particular, intersections of columns below are
semantic conjunctions, not automatically jointly executable tests on
one destructive specimen.

\begin{definition}[Generated experiments]
A generated experiment at $i$ is a labelled pair $(p,e)$ with
$p:i\to j$ and $e\in O_j$, evaluated by
\[
 C_{(p,e)}=X(p)^{-1}(C_e).
\]
Write $G_i$ for the family of its distinct response columns, while
retaining the pairs as native witness annotations. Define
\[
 M_i=\Meet(G_i)
 =\left\{\bigcap\mathcal U:\mathcal U\subseteq G_i\right\},
 \qquad
 \cl_i(A)=\bigcap\{C\in G_i:A\subseteq C\}.
\]
Empty intersections are $X_i$. Define $x\sim_i y$ when all columns of
$G_i$ have the same membership value at $x$ and $y$, and put
$Q_i=X_i/{\sim_i}$.
\end{definition}
The family $G$ is the least family containing the seeds and stable
under preimage along admitted programs. Precomposition acts on
\emph{labels} by $(q,e)\mapsto(qp,e)$, so the semantic construction has
actual native witnesses, not selected endpoint arrows.

\section{Generated state and concept actions}
\begin{theorem}[The source-derived operational doctrine]
For every $p:i\to j$:
\begin{enumerate}
\item $X(p)$ descends uniquely to $\bar p:Q_i\to Q_j$. These maps
make $Q$ a functor, with quotient maps $q_i:X_i\twoheadrightarrow Q_i$
natural in $i$.
\item Preimage restricts to a complete-meet map
$p^*:M_j\to M_i$. Its left adjoint is
\[
 p_!(A)=\cl_j(X(p)[A]).
\]
It preserves all joins and satisfies
$p_!\gamma_i(x)=\gamma_j(X(p)x)$, where
$\gamma_i(x)=\cl_i(\{x\})$.
\item The maps $p_!$ preserve identities and composition. They are the
unique complete-join actions agreeing with native state operations on
all marked object concepts $\gamma_i(x)$.
\end{enumerate}
\end{theorem}
\begin{proof}
A target generated column is $X(r)^{-1}(C_e)$, whose preimage under
$X(p)$ is $X(rp)^{-1}(C_e)\in G_i$. Consequently source-row equality
implies target-row equality. Surjectivity of $q_i$ gives uniqueness and
functoriality of the quotient operations.

Every extent is an intersection of generated columns; preimage preserves
arbitrary intersections, including the empty intersection. It therefore
restricts as claimed. For extents $A\in M_i,B\in M_j$,
\[
 \cl_j(X(p)[A])\subseteq B
 \quad\Longleftrightarrow\quad
 A\subseteq X(p)^{-1}(B),
\]
which proves the adjunction. Since $X(p)^{-1}(\cl_j(X(p)[U]))$ is a
closed set containing $U$, it contains $\cl_i(U)$. Thus
\begin{equation}
 \cl_j(X(p)[\cl_i(U)])=\cl_j(X(p)[U]).\tag{CN}
\end{equation}
This yields the marking equation and the composition law.
Finally $A=\bigvee_{x\in A}\gamma_i(x)$, with joins given by closed
unions. Every complete-join map agreeing on these generators is forced.
\end{proof}

For a one-object category this is a deterministic action of a monoid,
and generated tests are continuations. For a DFA with accepting set $F$,
$G=\{u_w^{-1}(F):w\in\Sigma^*\}$. The quotient is the predictive
row quotient; after retaining only reachable states it is the usual
minimal DFA. The lattice is its FCA/complete-meet completion, not the
state set itself. These fixed-interface facts are established automata,
FCA and Chu-space theory and were already proved in the preceding
\emph{common generation} manuscript.

\section{Growth forces a mixed-variance comparison}
Let $J:\mathcal C\hookrightarrow\mathcal D$ extend the admitted
program category, possibly by adding interface objects as well as
programs. Let $Y:\mathcal D\to\mathsf{Set}$ extend the source semantics,
with $YJ=X$. Old seed observations are retained at old interfaces;
new seeds may be supplied at any interface of $\mathcal D$. Decorate
new generated columns, lattices, closures and quotients by $+$.
At every old interface,
\[
 G_i\subseteq G_i^+,\qquad M_i\subseteq M_i^+.
\]
The new experiments can pass through newly introduced interfaces.
They need not be old experiments with renamed labels.

\begin{theorem}[Two directions, and strict backward coherence]
At each old interface the extension induces:
\begin{enumerate}
\item a unique surjection $\rho_i^Q:Q_i^+\twoheadrightarrow Q_i$
with $\rho_i^Qq_i^+=q_i$;
\item a complete-meet inclusion $j_i:M_i\hookrightarrow M_i^+$,
whose left adjoint is the surjective complete-join map
\[
 \rho_i^M:M_i^+\twoheadrightarrow M_i,
 \qquad \rho_i^M(A)=\cl_i(A);
\]
\item strict naturality for every \emph{shared} native program $p:i\to j$:
\begin{align*}
 \rho_j^Q\bar p^+&=\bar p\rho_i^Q,\\
 \rho_j^M p_!^+&=p_!\rho_i^M,\\
 \rho_i^M\gamma_i^+(x)&=\gamma_i(x);
\end{align*}
\item forward lax naturality, in the opposite direction,
\begin{align*}
 p_!^+j_i&\leq j_jp_!,\\
 \gamma_i^+(x)&\subseteq j_i\gamma_i(x).
\end{align*}
All backward comparisons compose strictly under successive extensions.
The forward inequalities are compatible with identities, composition
of shared programs, and successive extensions.
\end{enumerate}
\end{theorem}
\begin{proof}
Equality on more tests implies equality on fewer, proving (1).
For $A\in M_i^+$ and $B\in M_i$,
$\cl_i(A)\subseteq B$ iff $A\subseteq B$, so
$\rho_i^M\dashv j_i$. It is surjective because it fixes $M_i$.
The state square follows by evaluating it on all $x\in X_i$.
On right adjoints there is an equality
$j_i p^*=p^{*+}j_j$, since both sides are the same set-theoretic
inverse image. Taking left adjoints gives the second square. The
marking square follows from
$\cl_i\cl_i^+=\cl_i$.

New closure is smaller: $\cl_j^+(U)\subseteq\cl_j(U)$.
Apply this to $U=X(p)[A]$ to obtain the forward action inequality;
the marking inequality is the singleton case. For two shared programs,
monotonicity gives
\[
 q_!^+p_!^+j_i\leq q_!^+j_jp_!
                  \leq j_kq_!p_!.
\]
These are precisely the unique order comparisons between their
specified endpoints, so they respect the strict program composition
law. For successive growth stages $0\subseteq1\subseteq2$,
$\cl_0\cl_1=\cl_0$ and the inclusions compose literally;
the quotient comparison is unique on source-state representatives.
This proves all stated coherences.
\end{proof}
A newly added program need not act on the old quotient or lattice at all.
The theorem asserts naturality for shared operations, not an impossible
old action for every newly executable operation.

\begin{theorem}[Least simultaneous operational repair]
The generated extent family is the least Moore family containing its
seed columns for which every admitted source program has a
complete-join action agreeing with that program on marked object
concepts. Relative to an old interface, it is the least such extension
of the old extent family and all new seed columns. Independently, its
row equivalence is the largest program-stable equivalence refining the
seed distinctions, hence the coarsest adequate deterministic quotient.
\end{theorem}
\begin{proof}
For any Moore families $N_i$, a source map $f:X_i\to X_j$ has a
marking-compatible complete-join extension iff
\[
 f^{-1}(B)\in N_i\quad\text{for every }B\in N_j.\tag{E}
\]
For necessity, a join extension implies
$f[\cl_{N_i}(U)]\subseteq\cl_{N_j}(f[U])$, because
$\cl_{N_i}(U)$ is the join of the generators of $U$. Set
$U=f^{-1}(B)$: its closure still maps into $B$, so $U$ is closed.
For sufficiency, condition (E) proves (CN) and the adjunction exactly
as above. If every native primitive satisfies (E), induction on
programs forces every generated seed preimage into $N_i$; closure
under intersections forces $M_i\subseteq N_i$. Conversely $M$ itself
satisfies (E).

If an equivalence respects seeds and is stable under every native
program, equivalent states have equal membership in every generated
column, so it refines the constructed row equivalence. The latter is
itself stable by the first theorem. Neither part asserts that the
coarsest row quotient already supplies the least concept repair.
\end{proof}

\section{Three independent exactness conditions}
For clarity take a finite carrier $X$, old columns $F\subseteq\Pow(X)$,
new seed columns $N$, and the \emph{entire resulting} native generator
set $A$, including old operations acting on newly introduced tests.
Old columns can already have been generated by the old operation set,
but that does not make old operations redundant on $N$. Write
\[
 F^+=\Gen_A(F\cup N),\quad M=\Meet(F),\quad M^+=\Meet(F^+).
\]
The object nerve $K_F$ has nonempty faces the nonempty subsets of an
actual column of $F$. Include the empty simplex separately.

\begin{proposition}[Exact repair tests]
\begin{enumerate}
\item $Q_{F^+}\to Q_F$ is bijective iff every new seed is constant on
old row classes and every $a\in A$ respects old row equality.
\item $M=M^+$ iff $N\subseteq M$ and
$a^{-1}(B)\in M$ for all $a\in A,B\in M$. It suffices to check
$B\in F$.
\item $K_F=K_{F^+}$ iff every nonempty new seed and every nonempty
$a^{-1}(C)$, for $a\in A,C\in F$, is contained in an old column.
\end{enumerate}
The second condition implies the first. Neither forces the third.
The third need not imply either of the others.
\end{proposition}
\begin{proof}
For (1), old-class constancy and operation stability propagate through
all words; conversely an unchanged row quotient is stable for all
new operations and observes all new seeds. For (2), inverse image
preserves intersections, so the stated generator check gives closure
of $M$ under all $A$-words. Minimality then gives equality. Necessity
follows from the stability of $M^+$.

For (3), a nonempty column is itself a face because $X$ is finite.
Thus unchanged complexes force each indicated column to be covered.
Conversely let $D(F)$ be the family of subsets covered by old columns,
together with the empty set. The one-step covering condition implies
$a^{-1}(S)\in D(F)$ whenever $S\in D(F)$, by choosing a column
$C\supseteq S$ and using
$a^{-1}(S)\subseteq a^{-1}(C)$. The extra seeds are in $D(F)$;
induction gives $F^+\subseteq D(F)$ and therefore equal nerves.

Every old extent is a union of old row classes, proving the implication.
The CNOT example below separates (1) from (2), and its reset stage
separates (2) from (3). Conversely, an old universal column already
makes the nerve a simplex; adding arbitrary separating columns can
change both rows and the lattice without changing that nerve.
\end{proof}
For infinite carriers, the nerve criterion must say that every finite
subset of each added column is old-covered. A single whole-column
cover need not follow. The exactness assertion here concerns equality
of simplicial complexes, not an equivalence of their homotopy types.

\section{Supported topology and the actual partial-action domains}
Assume the state carriers in this section are finite. Define
\[
 \mathcal C_F=\{B\in\Meet(F):B\ne\varnothing,
                     \exists C\in F\ (B\subseteq C)\}.
\]
This is the supported concept poset. Closure gives a retraction
$r_F:\mathcal F(K_F)\to\mathcal C_F$, $U\mapsto\cl_F(U)$,
where $\mathcal F$ denotes nonempty faces. Its inclusion $s_F$ satisfies
$r_Fs_F=1$ and $1\leq s_Fr_F$, hence gives an order-homotopy
equivalence. The empty extent and the absence of a common column are
handled by the support condition, not by automatically deleting both
lattice extrema.

For a shared native program $p:i\to j$ at stage $k$, define
\begin{align*}
 K_{p,k}&=\{\varnothing\}\cup
 \{\varnothing\ne U\subseteq X_i:X(p)[U]\in K_{j,k}\},\\
 \mathcal C_{p,k}&=\{B\in M_{i,k}:B\ne\varnothing,
                         p_{!,k}(B)\in\mathcal C_{j,k}\}.
\end{align*}
Backward stability of columns makes $K_{p,k}$ a subcomplex of $K_{i,k}$.
These are the maximal domains on which the state action and concept
action have supported values. They need not be full domains.

\begin{theorem}[Growth of the supported bridge]
For growth from stage $0$ to stage $1$:
\begin{enumerate}
\item $K_{i,0}\subseteq K_{i,1}$ and
$\mathcal C_{i,0}\subseteq\mathcal C_{i,1}$.
The closure comparisons satisfy
\[
 r_{i,1}(U)\subseteq r_{i,0}(U)
 \qquad(U\in\mathcal F(K_{i,0})),
\]
inside the new supported concept poset. This inequality is a specified
order homotopy between the two forward comparison routes.
\item There are inclusions $K_{p,0}\subseteq K_{p,1}$ and
$\mathcal C_{p,0}\subseteq\mathcal C_{p,1}$. They give commuting
maps of the simplicial action spans
\[
 K_{i,k}\longleftarrow K_{p,k}\xrightarrow{X(p)}K_{j,k}.
\]
For the concept action spans, the left inclusion squares commute and
the right squares satisfy $p_{!,1}j\leq jp_{!,0}$. At each fixed
stage, the closure maps commute strictly with both arrows of the
action spans, including their domains.
\item For composable shared programs $p,q$,
\[
 K_{qp,k}=\{U\in K_{p,k}:X(p)[U]\in K_{q,k}\},
\]
with the identical formula for concept domains using $p_{!,k}$.
The stage-growth inequalities are compatible with these compositions.
\item The backward lattice projection is supported precisely on
\[
 D=\{B\in\mathcal C_{i,1}:\exists C\in G_{i,0},\ B\subseteq C\}.
\]
On $D$, $\rho^M=\cl_{i,0}$ maps to $\mathcal C_{i,0}$ and
retracts its inclusion. In general it is not defined as a supported
map on all of $\mathcal C_{i,1}$.
\end{enumerate}
\end{theorem}
\begin{proof}
An old column remains a new column. An old extent remains closed in
the finer Moore family and retains its support, proving (1). New
closure is smaller and has the same common-column property, so both
ends of the inequality lie in the new supported poset. For monotone
maps $u\leq v$, the map on the product with $\{0<1\}$ selecting
$u$ and $v$ is monotone; its realization is the claimed homotopy.

If $X(p)[U]$ is old-covered it is new-covered; this proves the
simplicial domain inclusion. An old closed $B$ remains new closed and
$X(p)[B]$ has its old cover, giving the concept-domain inclusion.
The simplicial squares commute because their vertex maps are the same
$X(p)$. The concept inequality was proved above. If a face $U$ has
$p$-image in a column $C$, then $U\subseteq X(p)^{-1}(C)$, an extent.
Consequently $\cl_i(U)$ remains in the same action domain.
Equation (CN) proves strict commutation with the action arrow; the
inclusion arrow is immediate. These domain closure maps are again
retractions with $1\leq sr$.

For (3), if $X(qp)[U]$ has a common target column, its preimage under
$q$ is a common column for $X(p)[U]$, and its further preimage is one
for $U$. This proves both directions of the domain formula. Use the
same common-column condition and (CN) for concepts. The lax coherence
is the chain of monotone inequalities already proved.

Finally $\cl_0(B)$ has an old common column iff $B$ does. This is
exactly the displayed domain. On it, $B\subseteq\cl_0(B)$ and
$\cl_0$ fixes old extents. Thus it retracts the inclusion of the old
supported core, with the stated order homotopy.
\end{proof}
\begin{remark}[Do not erase domain growth]
The result is a map of action spans, not automatically a natural
transformation of partial maps with equal domains. In fact
\[
 K_{p,0}\subseteq K_{i,0}\cap K_{p,1}
\]
can be strict: a newly available test can support the image of an old
simplex that previously lost all common tests. There is no global
homology endomorphism for a genuinely partial action. Likewise, the
forward concept inclusion generally fails to preserve the state
markings exactly; backward projection preserves them exactly.
\end{remark}

\section{Deriving when wiring and completion commute}
Let $P\subseteq X\times Y$ be the \emph{actual} compatible joint-state
carrier, with projections $\pi_X,\pi_Y$. Component seed columns are
$E_X,E_Y$, and their native completions are $G_X,G_Y$.
Let $T$ be specified joint native generators $t:P\to P$; these may
be interacting operations, not products. Wire only the seeds first:
\[
 F_0=\{\pi_X^{-1}(C):C\in E_X\}
       \cup\{\pi_Y^{-1}(D):D\in E_Y\}.
\]
Completing separately before wiring gives
\[
 F_{\mathrm{loc}}=\{\pi_X^{-1}(C):C\in G_X\}
       \cup\{\pi_Y^{-1}(D):D\in G_Y\}.
\]
Completing the joint system gives $F_{\mathrm{joint}}=\Gen_T(F_0)$.
The relation between these families must be proved from native
capabilities, not inferred from a prospective product of quotients.

\begin{definition}[Component experiment liftability]
For each native component program $u$ on $X$ there is an admitted
joint word $\widetilde u$ with
$\pi_X\widetilde u=u\pi_X$; likewise for $Y$.
The equality concerns actual source transformations on all of $P$.
The lifting words and their resource admissibility are supplied native
certificates. They need not leave the other component unchanged.
\end{definition}
\begin{theorem}[Exact operational completion--wiring comparison]
Assume component experiment liftability. Then
\[
 F_{\mathrm{loc}}\subseteq F_{\mathrm{joint}}
 =\Gen_T(F_{\mathrm{loc}}).
\]
Consequently joint completion has the forced comparison directions
of the growth theorem. In particular put
\[
 L=\im\bigl(P\xrightarrow{(q_X\pi_X,q_Y\pi_Y)}Q_X\times Q_Y\bigr).
\]
The local row quotient is $L$, not an assumed entire product. Then:
\begin{enumerate}
\item The map $Q_{\mathrm{joint}}\twoheadrightarrow L$ is an
isomorphism iff every joint primitive respects the fibers of the
displayed profile map.
\item The local extent family is
\[
 M_{\mathrm{loc}}=
 \{P\cap(A\times B):A\in M_X,\ B\in M_Y\}.
\]
Full FCA completion commutes with wiring iff every joint primitive
pulls every such closed rectangle back to another such rectangle.
It suffices to check the two kinds of generating local columns.
\item The local and joint object nerves are equal iff every nonempty
one-step inverse image $t^{-1}(C)$ of a local column is contained in
a local column. Here the carriers are finite.
\end{enumerate}
If every joint primitive is componentwise native,
$t(x,y)=(u(x),v(y))$ on $P$, then the literal column families themselves
are equal. Combined with liftability this supplies a sufficient
source-level commutation theorem for all the characterizations.
\end{theorem}
\begin{proof}
For a generated component experiment $u^{-1}(C)$, its local cylinder
is $\widetilde u^{-1}(\pi_X^{-1}(C))$, hence belongs to the joint
completion. This proves the first inclusion. Since
$F_0\subseteq F_{\mathrm{loc}}\subseteq\Gen_T(F_0)$, applying the
idempotent generation operator gives the asserted equality.

Two joint states agree on all local tests exactly when their two
component quotient profiles coincide. This proves the description of
$L$ and (1) by the repair criterion. Intersecting any family of local
cylinders yields $P\cap(A\times B)$, and every such rectangle is an
intersection of local cylinders. This proves the displayed formula;
(2) and (3) now follow from the respective repair criteria. Finally a
componentwise native primitive pulls an $X$-column back along $u$ and
a $Y$-column back along $v$. Both are again local generated columns.
Thus $F_{\mathrm{loc}}$ is already $T$-stable, proving literal equality.
\end{proof}
The liftability assumption is substantive with shared stock, ownership,
feedback, or compatibility constraints. If it fails, even the displayed
inclusion of local experiments in joint experiments is not automatic.
Separate component characterizations may posit experiments that cannot
be executed together on the actual compatible carrier. That is a native
source obstruction, not a failure repaired by declaring a quotient map.

\section{A native interaction that repairs FCA without refining states}
Take $X=Y=\{0,1\}$ and $P=X\times Y$, ordered $00,01,10,11$.
The component observations test $0$ and $1$; their initial native
operation categories contain only identities. The four local columns
are
\[
 x=0,\quad x=1,\quad y=0,\quad y=1.
\]
There are four distinct rows. The full extent family has ten sets:
$\varnothing$, four singleton sets, these four two-element columns,
and $P$. The object nerve is a square.

Now admit the actual reversible native operation
\[
 T(x,y)=(x,x\oplus y).
\]
Since $T^2=1$, all generated columns are exactly the four local columns
and
\[
 x\oplus y=0,\qquad x\oplus y=1.
\]
No additional column is concealed in a longer word. These six columns
are all six two-element subsets of $P$. The row quotient still has
four states, but the full extent family now has twelve sets: empty,
all four singletons, all six pairs, and $P$. The nerve is the graph
$K_4$, with no triangular faces, so
\[
 \beta_1(K_{\mathrm{local}})=1,
 \qquad \beta_1(K_{\mathrm{new}})=6-4+1=3.
\]
These are Betti numbers over $\mathbb F_2$, and also over every field.

The FCA obstruction is explicit. For $U=\{00,11\}$,
\[
 \cl_{\mathrm{local}}(U)=P,
 \qquad
 T[U]=\{00,10\},
 \qquad
 \cl_{\mathrm{local}}(T[U])=\{00,10\}.
\]
Thus $T[\cl_{\mathrm{local}}(U)]\not\subseteq
\cl_{\mathrm{local}}(T[U])$. No complete-join action on the old
concept lattice can agree with CNOT on its four marked object states.
The old state quotient is already discrete, so CNOT acts on it without
any repair. The twelve-element lattice is the least operational FCA
repair. This is a source-derived separation between the state and
concept versions of completion--wiring commutation.

Next admit the native reset $R(x,y)=(x,0)$. Closing under all mixed
$T,R$ words gives exactly
\[
 \{x=0,x=1,y=0,y=1,x\oplus y=0,x\oplus y=1,\varnothing,P\}.
\]
Indeed $R^{-1}(y=0)=P$, $R^{-1}(y=1)=\varnothing$;
$R^{-1}(x\oplus y=b)=\{x=b\}$;
$R^{-1}(x=b)=\{x=b\}$; and $T$ permutes the six pair-columns
and fixes $\varnothing,P$. Hence this eight-column family is closed
under both primitive inverse images and contains all one-step new
columns, proving exact all-word closure. The row quotient and extent
family remain four and twelve respectively: $\varnothing$ and $P$
were already extents. But the actual universal test $P$ makes the
nerve a full $3$-simplex, so $\beta_1$ drops from $3$ to $0$.

The prior common-generation manuscript already showed that a universal
column can change support topology without changing an unmarked full
lattice. The added content here is not that elementary phenomenon.
It is the independently specified native interaction, its complete
mixed-word experiment generation, the minimal FCA repair, and the
exact source-level separation between the three commutation criteria.

\section{An automata change and noncommuting repairs}
A reachable automata example uses states $\{0,1,2\}$, initial state $2$,
accepting set $\{0\}$, and the original letter $a$ with
\[
 a(0)=a(1)=a(2)=0.
\]
Its generated columns are $\{0\}$ and the full carrier. Its predictive
classes are $\{0\}$ and $\{1,2\}$, and both quotient states are
reachable. Add the native letter $b$ with
\[
 b(0)=0,\qquad b(1)=0,\qquad b(2)=1.
\]
The sole additional column is $b^{-1}(\{0\})=\{0,1\}$;
all three source states are now reachable from $2$, by $\epsilon,b,bb$.
States $1,2$ are separated by the executable continuation $b$.
The minimal recognizer therefore grows from two to three states, and
the extent chain from two to three elements. The nerve remains the
full simplex because the old experiment $a$ already accepts all source
states. Thus native automata refinement need not change topology.

There is a separate obstruction to performing each interface repair
only once. On $\{0,1,2,3\}$, with seed $\{1\}$, take
\[
 a=(0,0,0,1),\qquad b=(0,0,3,0).
\]
First closing the seed under $b$ and then under $a$ gives exactly
$\varnothing,\{1\},\{3\}$. It identifies $0$ and $2$.
Full mixed-word closure additionally contains $\{2\}$, witnessed by
executing $b$ then $a$ and then the seed observation. The complete
family is exactly $\varnothing,\{1\},\{2\},\{3\}$, since inverse
images of these four sets remain among them. The once-completed
interface is not stable under $b$ after $a$ has introduced $\{3\}$.

\begin{proposition}[Exact one-pass condition]
Let $S_A,S_B$ be operational-saturation closure operators on Moore
families, for two native generator sets. Then
\[
 S_{A\cup B}(M)=S_A S_B(M)
\]
iff $S_A S_B(M)$ is also $B$-stable. In general the union completion
is the least common fixed point above $M$, obtained by alternating
saturation; on a finite carrier this process terminates. If the two
native generated submonoids commute setwise, the two saturation
operators commute and one pass in either order suffices.
\end{proposition}
\begin{proof}
The one-pass family is $A$-stable and contains $M$. It is the full
completion exactly when it is also $B$-stable, by leastness and the
inclusion $S_A S_B(M)\subseteq S_{A\cup B}(M)$.
Every common fixed point contains every alternating iterate, and their
union followed by meet closure at limit stages is the least common
fixed point. On a finite carrier only finitely many subsets can be
added. If the generated submonoids commute setwise, every mixed word
can be rewritten as an $A$-block followed by a $B$-block and also in the
opposite order. Their inverse-image families, and hence their meet
closures, coincide.
\end{proof}

\section{Connecting a decomposition comparison to actual incidence}
An exactness theorem about labelled experimental decompositions does
not automatically identify any of the state or concept constructions.
The missing link is a response evaluation, and its variance must be
explicit. Here is the exact transfer statement needed to connect a
Kan-extension comparison to this operational bridge.

Fix an actual source-state carrier $P$. A labelled experiment family
with responses is an object $(E,\mathrm{ev})$ of the slice category
$\mathsf{Set}/\Pow(P)$, where $\mathrm{ev}(e)$ is the accepting
column of experiment $e$. A map in this slice preserves the entire
response column, not just the existence of an experiment.

\begin{proposition}[Response-cone transfer]
Suppose source and target decomposition diagrams carry compatible
response cones to the constant set $\Pow(P)$, and a canonical
comparison of their colimits
\[
 \Theta:E_{\mathrm{source}}\longrightarrow E_{\mathrm{target}}
\]
is induced by a functor of decomposition categories. Then
$\mathrm{ev}_{\mathrm{target}}\Theta
=\mathrm{ev}_{\mathrm{source}}$.
If $\Theta$ is bijective, it gives an isomorphism of labelled incidence
contexts. More weakly, equality of the two sets of evaluated columns
already identifies their row quotients, full extent families, supported
cores and object nerves. Surjectivity of $\Theta$ suffices for that
weaker conclusion. If the labelled comparison also commutes with
native program precomposition, the identifications preserve the
induced state and concept actions and the exact partial-action domains.
\end{proposition}
\begin{proof}
A compatible response cone is constant on the identifications defining
the colimit, so it descends uniquely to its evaluation map. Evaluating
a source representative before or after the comparison uses the same
cone value, proving the displayed equality. A bijective label map with
this equality preserves incidence in both directions. All four
extensional constructions depend only on the family of evaluated
columns; this proves the weaker statement. Finally precomposition
compatibility identifies the columns witnessing action domains and the
source state maps, so the formulas of the earlier theorems agree.
\end{proof}
For independently executable parallel tests, a possible evaluation is
$\pi_X^{-1}(C_e)\cap\pi_Y^{-1}(C_f)$, with actual parallel execution
witnesses. It is a valid response cone only when the admitted
refinements preserve these columns; independent execution cannot be
inferred from a formal product of evidence sets. Other experiment
composition operations need their own native evaluation.

A cofinality theorem for labelled decompositions can therefore imply
an incidence theorem through this explicit response contract. Its
failure need not change any of the four extensional characterizations:
missing or duplicated labels can evaluate to already represented
columns. Likewise, the CNOT interaction is not automatically a Day
convolution of tests. It is an actual state transformation whose
preimages generate additional response columns. Identifying these
constructions without the response cone and a native composition
interpretation would be only an analogy, not a theorem.

\section{Verification, predecessors, and exact increment}
The executable companion \texttt{check\_evolving\_incidence.py} uses
bitmask columns and exact finite-state preimages. It checks complete
word closure, retains an explicit word for each column, computes meet
closures and row classes, and computes simplicial Betti numbers by
Gaussian elimination over $\mathbb F_2$. The CNOT/reset calculation
includes empty columns and all mixed words, rather than stopping after
one inverse-image round. Exhaustive tests cover every distinct-column
context on at most three objects and every total unary operation:
6,980 contexts-with-operation, 20,940 exactness comparisons, and
173,796 shared-operation/refinement-square checks. A separate script
records the reachable three-state automata example. These finite checks
are sanity checks for the proofs, not a replacement for them.

The individual ingredients are established: inverse-image predicate
transformers, Moore-family completion, deterministic congruence
refinement, FCA adjunctions, Chu transport, and the supported
FCA--Dowker comparison. Angluin's observation-table work and subsequent
categorical automata-learning frameworks already make the addition of
distinguishing continuations central. Formal-context morphism theory
already distinguishes lattice adjunctions from other map classes.
Functorial Dowker and nerve theorems already give broad topological
comparison principles. No priority claim for those ingredients is made.

Within this project, \emph{common\_generation\_note.tex} supplies the
fixed acted-incidence bridge, its exact extension obstruction, and its
maximal partial topology. \emph{adaptive\_experiment\_formation.tex}
supplies finite transition-stable repair, its executable certificates,
and exact belief-update/native-policy transport.
\emph{contextual\_reduction\_wiring.tex} supplies the realizability-image
comparison for fixed contextual quotients and separates several notions
of adequacy. \emph{native\_operations\_across\_doctrines.tex} separates
semantic operation reconstruction from native realizability and coherent
history lifting. The present development does not supersede those
obligations.

The increment is their changing-interface organization: generated
experiments force opposite-direction strict quotient/FCA comparisons
and forward lax supported comparisons; common native operations and
partial domains obey explicitly proved coherence; operational completion
commutes with wiring exactly under the distinct state, rectangular
FCA, and support criteria; and a concrete native interaction separates
those criteria. The response-cone proposition states exactly how a
separate categorical theorem about changed decomposition interfaces
can enter this same semantics. Finite observations still do not prove
an unrestricted source model, and extensional agreement still does not
recover erased execution witnesses or residual resources.

\end{document}
